\documentclass[11pt,a4paper]{article}
\usepackage[T1]{fontenc}
\usepackage[utf8]{inputenc}
\usepackage[francais]{babel}
\usepackage{color}
\usepackage{graphicx}
\title{
	\textcolor{blue}{Le Mans Université}\\
	Licence Informatique 2ème année\\
	Module 174UP02 Rapport de Projet\\
	\textbf{For The King}\\
}
\author{Durand Leo, Tsamarayev Massoud, Mohamed Saandi}
\date{\today}
\begin{document}

\maketitle
\newpage

\tableofcontents
\newpage

\section{Introduction}

\textit{Cette introduction présentera le sujet qui sera traité et le travail avec une présentation du plan adopté}

Ce document présente le projet de jeu For The King réalisé dans le cadre de la formation de L2 informatique de l'université du Mans pendant la période de janvier à avril 2026.
Ce projet a été développé en langage C avec la librairie SDL.
Ce jeu fait ceci cela...


\section{Conception}

blablacar

\subsection{Présentation du jeu}

For the king est un jeu tour par tour, basé sur de l'aléatoire.

\subsection{Fonctionnalités du jeu}

fonc

\section{Organisation du projet}

1234

\subsection{Planning prévisionnel}

ab

\subsection{Répartition des tâches}

Massoud : génération de l'eau, affichage des hexagones, mécaniques du combat, menus, graphismes, système de niveau, affichage des combats, placement des monstres, quêtes
Leo : placement du personnage sur la carte, mécaniques de déplacement, inventaire, système de sauvegarde, affichage de la carte en SDL, affichage des combats, tests unitaires, boss final
Saandi : documentation Doxygen, génération des biomes, affichage des informations lors des interactions en SDL, Gantt prévisionnel.
Toute l'équipe : création des monstres et leur attributs, création des personnages et leur attributs, carte

\subsection{Outils de travail}

ef

\section{Développement}

gh

\subsection{Génération aléatoire de la carte}

ij

\subsection{Affichage du jeu}

12

\subsection{Gestion des combats}

34

\subsection{Gestion du menu}

56

\subsection{Gestion du personnage}

78

\subsection{Gestion de la sauvegarde}

90

\section{Bilan et résultats}

titi

\section{Bibliographie}

toto

\section{Annexe}

aaaa

\end{document}
